\documentclass[11pt,a4paper]{article}
\usepackage[utf8]{inputenc}
\usepackage[T1]{fontenc}
\usepackage[brazil]{babel}
\usepackage[a4paper,margin=2.3cm]{geometry}
\usepackage{enumitem}
\usepackage{hyperref}
\usepackage{xcolor}
\setlength{\parindent}{0pt}
\setlength{\parskip}{5pt}
\sloppy
\setlength{\emergencystretch}{3em}
\hypersetup{colorlinks=true,urlcolor=blue,linkcolor=black}
\newcommand{\cabecalho}[2]{\begin{center}{\small \textbf{Sprint 2 -- #1}}\\[1mm]\textit{Inventário de Ativos e Vulnerabilidades}\\Matheus Moreli Durão Silva -- 12621CBS217\\Cibersegurança -- Bacharelado -- Noturno\end{center}\hrule\vspace{4mm}}
\newcommand{\bloco}[1]{\par\vspace{2mm}\textbf{#1}\par}
\newcommand{\codigo}[1]{\texttt{#1}}

\begin{document}\cabecalho{S2\_17 -- Colaboração no GitHub}{}
\bloco{Objetivo} Colaboração não é apenas apertar ``merge''. A Pull Request deve deixar claro o que mudou, como testar e qual card/requisito ela atende, para que outra pessoa consiga revisar sem adivinhar.
\bloco{Modelo de descrição} ``Card: S2\_08. Mudança: cadastro de ativo recebeu criticidade e validação de ID único. Teste: executar python -m unittest -v test\_sprint2.py. Resultado esperado: todos os testes OK.'' Não colocar senhas, tokens, dados pessoais ou chaves em código, commit, comentário ou captura de tela.
\bloco{Revisão real} O checklist pede aprovação ou comentário de um colega e resposta aos comentários antes do merge. Essa parte não pode ser simulada: é necessário convidar uma pessoa real ao repositório e registrar a revisão na Pull Request.
\bloco{Checklist} PR referencia card, explica alteração/teste, não expõe segredo, recebe revisão real e responde comentários antes da integração.
\end{document}
